
\section{Linear Theory}
\sectionframe{2}{Linear Theory: Minimum Variance Linear Estimator for Stochastic Processes}

%--------------------------
\myframe{Considerations}{%

\begin{tabular}{clcp{7cm}}
1. & Linear Estimators & $\Longrightarrow$ & depending on first and second
                                        statistics, sufficient to find
                                        estimators minimizing squared measures
                                        of the estimation errors, e.g.,
                                        variance.\\
 & & & \\
\pause
2. & Processes & $\Longrightarrow$ & so far we considered $\x$ and $\y$ of fixed
                                known dimensions. We now turn to the more
                                interesting (and more useful) setting of
                                \tq{continuous} flow of new measurements
                                $\{\y(t)\}$ of a (most often) time-varying
                                signal $\{\x(t)\}$. In this scenario, the time
                                interval $I$, during which data/measurements
                                $\{\y(t)\}_{t\in I}$ flow, increases. Hence, it
                                appears clearly how a suitable algorithmic
                                structure to process an increasing amount of
                                data becomes indeed fundamentally important.
\end{tabular}
}


%--------------------------
\myframe{Some more notation related to Hilber spaces...}{%
    \begin{itemize}
        \item When dealing with (random) processes $\y(t)$ (collection of
        variables) over time, i.e., $\{\y_k(t), t\in\Z, k=\until{m}\}$, the
        Hilbert space $\H(\y)\subset\H$ generated by the (second order) process
        can be defined as

        \begin{align*}
            \tilde{\H}(\y) & = \left\{\sum a_{ki}\y_k(t_i); a_{ki}\in\R, t_i\in\Z, k=\until{m}\right\}\\
            \H(\y) &= \clspn{\y(t), t\in\Z} = \cls{\tilde{\H}(\y)}
        \end{align*}
        
        with $\cls{\cdot}$ denoting the formal operation of closure that is
        completing the corresponding space with the limits of all its Cauchy
        sequences which, in the case of infinite dimensional space, could not
        own to the original space.\footnote{We will take these technical aspects
        for granted.}

        \pause
        \item Considering past and present time $s\in(-\infty, t]$, we denote
        $$
        \Hpast{t}(\y) :=\clspn{\y(s);\ s\leq t};\ \ \Hpast{s}(\y)\subseteq \Hpast{t}(\y), s\leq t;\ \ \lim_{t\to\infty}\Hpast{t}(\y)=\H(\y)
        $$

        \pause
        \item Considering present and future time $s\in[t,+\infty)$, we denote
        $$
        \Hfuture{t}(\y):=\clspn{\y(s);\ s\geq t};\ \ \Hfuture{s}(\y)\supseteq\Hfuture{t}(\y), s\leq t;\ \ \lim_{t\to-\infty}\Hfuture{t}
        (\y)=\H(\y)
        $$

    \end{itemize}
}

%--------------------------
\myframe{...and orthogonal projections}{%

Given $\H(\y)\subset\H$, $\forall \x\in\H$

\pause
$$
\x = \xhat + \tilde{\x}
$$

\pause
where $\xhat\in\H(\y)$ is the orthogonal projection (also denoted as
$\Elin[\x|\H(\y)]$) of $\x$ onto $\H(\y)$ and $\tilde{\x}\bot\H(\y)$

\vspace{.5cm}

\begin{itemize}
    \pause
    \item (idempotent) $\Elin\left[\Elin[\cdot|\H(\y)]|\H(\y)\right] = \Elin[\cdot|\H(\y)]$
    \pause
    \item if $\H(\y_1)\subset\H(\y_2)$, $\Elin\left[\Elin[\cdot|\H(\y_2)]|\H(\y_1)\right] = \Elin[\cdot|\H(\y_1)]$
\end{itemize}

\pause
In particular if $s<t$

$$
\Elin\left[\Elin[\x|\Hpast{t}(\y)]|\Hpast{s}(\y)\right] = \Elin[\x|\Hpast{s}(\y)]
$$
}

%-------------------------------------------------------------------------------
\subsection{State Space models}

\myframe{State Space models}{%
Let $\{\y(t)\}$ an $m$-dimensional stochastic process defined on $[t_0,
\infty)$. 
\pause
We say it is a \textit{finite dimensional process} if we can express
it as

\begin{align*}
    \x(t+1) &= A\x(t) + B\n(t)\\
    \y(t) &= C\x(t) + D\n(t),\quad t\geq t_0,\ \x(t_0)=\x_0
\end{align*}

where

\begin{itemize}
    \pause
    \item $A,B,C,D$ are \textit{deterministic} matrices, possibly time varying
    but know\footnote{To avoid notation clutter, time dependence of the model
    matrices is omitted.}.
    
    \pause
    \item $\{\n(t)\}$ is a (orthonormal) \textit{white noise} process with zero
mean and unit variance ($\E[\n(t)]=0,\ \E[\n(t)\n(s)^T]=I\delta_{ts}$)
    
    \pause
    \item $\{\x(t)\}$ is a $n$-dimensional process, called \textit{state
    process}, with initial value $\x_0$ which is uncorrelated w.r.t. all present
    and future values of $\n$, i.e., $\E[x_0\n(t)^T]=0$, $\forall t\geq t_0$ and
    $\E[\x_0]=\mu_0$, $Var[\x_0]=\Sigma_0$ 
\end{itemize}

\vspace{.1cm}

\pause
\begin{alertblock}{}
Essentially a finite dimensional process can be described as the output of a
(time-varying) linear dynamical system driven by white noise.
\end{alertblock}
}

%--------------------------
\myframe{State Space models - the state process}{%
    \begin{block}{Markov property}
        $\{\x(t)\}$ is a weak (sometimes also referred to as second-oder) Markov process, i.e.,
        
        $$
        \Elin[\x(t)|\H_s(\x)] = \Elin[\x(t)|\x(s)]\quad \forall t\geq s
        $$
        
        being $\H_s(\x)=\Hpast{s}(\x)$ the Hilbert space generated by
        $\{\x(t_0),\ldots,\x(s)\}$. If $\{\n(t)\}$ and $\x_0$ are jointly
        Gaussian, $\{\x(t)\}$ is strongly Markov (i.e., replace $\Elin$ by $\E$,
        equivalently the entire distribution).
    \end{block}

    \pause

    \textbf{Sketch of proof:}
    \vspace{-.4cm}
    \begin{align*}
        &\x(t) = A^{t-t_0}\x_0 + \sum_{k=t_0}^{t-1}A^{t-1-k}B\n(k) = A^{t-s}\x(s) + \sum_{k=s}^{t-1}A^{t-1-k}B\n(k), \quad t\geq s\\
        %
        &\Rightarrow \{\x(t)\} \text{linear in}\ \x_0,\ \{\n(t_0),\ldots,\n(t-1)\}%
        \Rightarrow \x(t)\bot\Hfuture{t}(\n)%
        \Rightarrow \Hpast{s}(\x) \bot \Hfuture{s}(\n)\\
        %
        &\Rightarrow \Elin[\x(t)|\Hpast{s}(\x)] = \Elin[\x(t)|\x(s)] = A^{t-s}\x(s)
    \end{align*}
    For the Gaussian case, $\Elin\equiv\E$. Also the entire conditional
    probability depends on the past only through the conditional mean hence the
    property holds \tq{globally}. 
}

%--------------------------
\myframe{State Space models - the state process}{%
The same \tq{memory} property holds wrt the output process $\{\y(t)\}$

\vspace{-.4cm}

\begin{align*}
    &\Elin[\x(t)|\Hpast{s}(\x, \y)] = \Elin[\x(t)|\x(s)]\\
    &\Elin[\y(t)|\Hpast{s}(\x, \y)] = \Elin[\y(t)|\x(s)],\quad t\geq s
\end{align*}

being $\Hpast{s}(\x, \y)=\Hpast{s}(\x) \vee  \Hpast{s-1}(\y)$ vector sum of
$\Hpast{s}(\x)$ (past and present wrt to $s$) and $\Hpast{s-1}(\y)$ (strict past
wrt to $s$).

\pause
\begin{block}{Conditional Incorrelation (Orthogonality)}
    Define $\z(t) := [\x(t+1)^T, \y(t)^T]^T$,
    \begin{align*}
        &\Hpast{t}(\z) := \Hpast{t+1}(\x) \vee \Hpast{t}(\y),\quad
        \Hfuture{t}(\z) := \Hfuture{t+1}(\x) \vee \Hfuture{t}(\y)\\
        &\X_t := \spn{\x_1(t),\ldots,\x_n(t)}\quad \text{state space of $\x$ at t} 
    \end{align*}
    
    then

    $$
    \Hpast{t}(\z)\ \bot\ \Hfuture{t}(\z)\ |\ \X_t
    \ \Longrightarrow\
    \hat{\z}(t) = \Elin[\z(t)|\x(t)] = \begin{bmatrix}A \\ C\end{bmatrix}\x(t)
    $$
\end{block}
}

%--------------------------
\myframe{State Space models - moments}{%

From $\E[\n]=0$ and previous (conditional) incorrelation properties, the
evolution of mean and covariance of a linear finite dimensional system is
described by

\begin{align*}
    \mux(t+1) &= A\mux(t),\quad \mux(t_0)=\muzero\\
    \muy(t)  &= C\mux(t)\\
    \Varx(t,s) &=
        \begin{cases}
            A^{t-s}\Sigma(s),\quad &t\geq s\\
            \Sigma(t)(A^T)^{s-t},\quad &t\leq s
        \end{cases}\\
    \Vary(t,s) &=
        \begin{cases}
            CA^{t-s-1}G(s),\quad &t>s\\
            C\Sigma(t)C^T + DD^T,\quad &t=s\\
            G(t)^T(A^T)^{s-t-1}C^T,\quad &t<s
        \end{cases}\\
    G(s) &= A\Sigma(s)C^T + BD^T\\
\end{align*}

and $\Sigma(t)=\Var[\x(t)]$ is the solution of the discrete Lyapunov equation

$$
\Sigma(t+1) = A\Sigma(t)A^T + BB^T,\quad \Sigma(t_0)=\Varzero
$$
}

%--------------------------
\myframe{State Space models - stability and (asymptotic) stationarity}{%

Stability and (asymptotic) stationarity\footnote{Evolution of covariance
depends on $\tau=t-s$, $\Sigma(s,t)=\Sigma(\tau)$, for $t,s\geq
t_0$ and $t-t_0\to \infty$.} of the underlying processes play an important role
in the relation between the Wiener-Kolmogorov and Kalman theories.

\vspace{.4cm}

\pause
Assuming $A,B,C,D$ are constant, if $\rho(A)<1$ (asymptotically stable):
\begin{itemize}
    \pause
    \item $\mux(t+1)=A\mux(t) \to 0$
    
    \pause
    \item $\Sigma(t+1)=A\Sigma(t)A^T+BB^T \to \Sigma_{\infty} = \sum_{k=0}^\infty A^kBB^T(A^T)^k$ (unique and $\succ0$)
    
    \pause
    \item If $\x_0\sim(0,\Sigma_\infty)$, then $\{\x(t)\}, \{\y(t)\}$ are weakly stationary ($\forall t\geq t_0$) with
    $$
    \Vary = C\Sigma_\infty C^T + DD^T.
    $$
    % \item If $\rho(A)\ge 1$, in general second moments are not bounded and the
    % process is not asymptotically stationary.
\end{itemize}

\pause
\begin{alertblock}{}
Under the above conditions\footnote{For finite LTI systems, stationarity
translates into rationality of the transfer function and, in turn, ARMA/ARMAX
representation. Hence the possibility to practically compute the filter transfer
function.} the two approaches (Wiener-Kolmogorov and Kalman) are essentially
equivalent.
\end{alertblock}

}

%--------------------------
\myframe{State Space models - a common alternative notation}{%

\begin{align*}
    \x(t+1) &= A\x(t) + B\n(t),\quad \E[\n(t)]=0,\ \E[\n(t)\n(s)^T]=I\delta_{ts}\\
    \y(t) &= C\x(t) + D\n(t),\quad t\geq t_0,\ \x(t_0)=\x_0
\end{align*}

\pause

By defining $\v(t):=B\n(t)$ (state/process noise) and $\w(t):=D\n(t)$
(measurement/observation noise) and $\z(t)=\begin{bmatrix}\x(t)^T & \y(t-1)^T\end{bmatrix}^T$, 
we have

$$
    \z(t+1) =
    \begin{bmatrix}A \\ C\end{bmatrix}\z(t) + \begin{bmatrix}\v(t)\\\w(t)\end{bmatrix},\quad
    \E\begin{bmatrix}\v(t)\\\w(t)\end{bmatrix} = 0,\
    \Var\begin{bmatrix}\v(t)\\\w(t)\end{bmatrix} = \begin{bmatrix}Q & S \\ S^T & R \end{bmatrix}
$$

Often, in practical scenarios, $S=0$ i.e. uncorrelated process and measurement noise.
}



%--------------------------
\myframe{(Some) Realization techniques}{%

(With a change of notation) Consider a discrete time ARMAX (Auto-Regressive Moving Average Exhogenous) model,
extremely common to represent input/output relations:

\pause
$$
A(z^{-1})\,y_k = B(z^{-1})\,u_k + C(z^{-1})\,e_k,\qquad e_k \sim (0,\sigma^2)
$$
%
\pause
\vspace{-.5cm}
\begin{align*}
&A(z^{-1}) = 1 + a_1 z^{-1} + \cdots + a_{n_a} z^{-n_a},\\
&B(z^{-1}) = b_0 + b_1 z^{-1} + \cdots + b_{n_b} z^{-n_b},\\
&C(z^{-1}) = 1 + c_1 z^{-1} + \cdots + c_{n_c} z^{-n_c}.
\end{align*}

\pause
\begin{itemize}
    \item ARMA: $B=0$ or $u=0$
    \item ARX: $C=1$ 
\end{itemize}

\pause
\textbf{Assumptions (for simplicity):}
\begin{itemize}
    \item Scalar (input and output, SISO) processes. For MIMO processes,
    generalization using matrix polynomials
    \item $A(z^{-1})$ and $C(z^{-1})$ are already monic (highest order coeff. =1)
\end{itemize}
}


%--------------------------
\myframe{Realization: Controllable Canonical Form}{%

The ARMAX description in transfer functions terms can be written:
\[
y_k = \frac{B(z^{-1})}{A(z^{-1})}u_k + \frac{C(z^{-1})}{A(z^{-1})}e_k
\]
\pause
An equivalent state space form with state $x_k\in\R^{n_a}$ is
\[
x_{k+1} = A_c\,x_k + B_c\,u_k + B_c\,e_k,
\qquad
y_k = (C_u+C_e)\,x_k + D_u\,u_k + D_e\,e_k.
\]
\pause
\vspace{-.5cm}
\begin{align*}
    A_c &=
        \begin{bmatrix}
            -a_1 & -a_2 & \cdots & -a_{n-1} & -a_n\\
            1 & 0 & \cdots & 0 & 0\\
            0 & 1 & \cdots & 0 & 0\\
            \vdots & & \ddots & & \vdots\\
            0 & 0 & \cdots & 1 & 0
        \end{bmatrix},
        \quad
    B_c=
        \begin{bmatrix}
            1\\0\\ \vdots\\0
        \end{bmatrix}
        \quad
        D_u=b_0,\quad D_e=1\\
    C_u &= \begin{bmatrix}b_1-a_1 b_0 & \cdots & b_n-a_n b_0\end{bmatrix}\quad
    C_e = \begin{bmatrix}c_1-a_1 & \cdots & c_n-a_n\end{bmatrix}
\end{align*}

\pause
\begin{alertblock}{}
\textbf{Controllable form}: controllability matrix from $(A_c, B_c)$ has
full rank, i.e., the system is controllable (input drives the state everywhere
in finite time).    
\end{alertblock}

}

%-----------------------
\myframe{Realization: Controllable  Form - Example (part 1)}{

Consider the ARMAX model
%
$$
y_k = \frac{b_0 + b_1 z^{-1}}{1 + a_1 z^{-1} + a_2 z^{-2}}u_k + \frac{1 + c_1 z^{-1}}{1 + a_1 z^{-1} + a_2 z^{-2}}e_k
$$
%
where 
\begin{align*}
    &A(z^{-1}) = 1 + a_1 z^{-1} + a_2 z^{-2},\\
    &B(z^{-1}) = b_0 + b_1 z^{-1},\\
    &C(z^{-1}) = 1 + c_1 z^{-1}.
\end{align*}

\pause
Using the above formulas, the state space matrix are
% 
\begin{align*}    
    A_c &=
        \begin{bmatrix}
            -a_1 & -a_2\\
            1 & 0
        \end{bmatrix},\
    B_c=
        \begin{bmatrix}
            1 \\ 0
        \end{bmatrix},\ 
    D_u=b_0,\ D_e=1,\\
    C_u &= \begin{bmatrix}b_1-a_1 b_0 & -a_2b_0\end{bmatrix},\
    C_e = \begin{bmatrix}c_1-a_1 & -a_2\end{bmatrix}
\end{align*}
%
\pause
corresponding to the transfer function
$$
y_k = (C_u(zI-A_c)^{-1}B_c + D_u)u_k + (C_e(zI-A_c)^{-1}B_c + D_e)e_k 
$$
}


%-----------------------
\myframe{Realization: Controllable  Form - Example (part 2)}{

Explicit computation lead (for the first term, the second can be similarly checked)
%
\begin{align*}
    C_u(zI-A_c)^{-1}B_c + D_u
    &= \begin{bmatrix}b_1-a_1 b_0 & -a_2b_0\end{bmatrix} \begin{bmatrix}z+a_1 & a_2\\-1 & z \end{bmatrix}^{-1} \begin{bmatrix}1\\0\end{bmatrix} + b_0\\ 
    &= \begin{bmatrix}b_1-a_1 b_0 & -a_2b_0\end{bmatrix} \frac{1}{z^2+a_1 z+a_2}\begin{bmatrix}z & -a_2\\1 & z+a_1 \end{bmatrix} \begin{bmatrix}1\\0\end{bmatrix} + b_0\\
    &= \frac{(b_1-a_1 b_0)z - a_2b_0 + b_0(z^2+a_1 z+a_2)}{z^2+a_1 z+a_2}\\
    &= \frac{b_1z  + b_0z^2}{z^2+a_1 z+a_2}\\
    &= \frac{z}{z^2}\frac{b_1  + b_0z}{1 + a_1 z^{-1} + a_2z^{-2}}\\
    &= z^{-1}\frac{b_1  + b_0z}{1 + a_1 z^{-1} + a_2z^{-2}} = \frac{b_0 + b_1z^{-1}}{1 + a_1 z^{-1} + a_2z^{-2}}
\end{align*}

}


%--------------------------
\myframe{Realization: Observable Canonical Form (Dual)}{%

A dual (observable) state space form with state $x_k\in\R^{n_a}$ is
\[
x_{k+1} = A_o\,x_k + B_u\,u_k + B_e\,e_k,
\qquad
y_k = C_o\,x_k + D_u\,u_k + D_e\,e_k.
\]
%
\pause
where the \textbf{observable canonical matrices} are built as the dual of the
controllable companion form, i.e., $A_o=A_c^\top$ and $C_o=B_c^\top$.

\pause
\vspace{-.5cm}
\begin{align*}
    A_o &= A_c^\top =
        \begin{bmatrix}
            -a_1 & 1 & 0 & \cdots & 0\\
            -a_2 & 0 & 1 & \cdots & 0\\
            \vdots & \vdots & & \ddots & \vdots\\
            -a_{n-1} & 0 & 0 & \cdots & 1\\
            -a_n & 0 & 0 & \cdots & 0
        \end{bmatrix},
        \quad
    C_o = B_c^\top =
        \begin{bmatrix}
            1 & 0 & \cdots & 0
        \end{bmatrix}\\
    %
    B_u &:=
    \begin{bmatrix}
        b_1-a_1 b_0\\
        b_2-a_2 b_0\\
        \vdots\\
        b_n-a_n b_0
    \end{bmatrix},
    \quad
    B_e :=
    \begin{bmatrix}
        c_1-a_1\\
        c_2-a_2\\
        \vdots\\
        c_n-a_n
    \end{bmatrix}\quad D_u=b_0,\quad D_e=1
\end{align*}

\pause
\medskip
\begin{alertblock}{}
\textbf{Observable form}: observability matrix from $(A_o, C_o)$ has
full rank, i.e., the state can be reconstructed from a finite window of outputs.
\end{alertblock}
}


%--------------------------
\myframe{Realization: Innovation (Kalman) Form}{%

An \textbf{innovation-form} (a.k.a. \textbf{Kalman form}) state-space model is
\[
x_{k+1} = A\,x_k + B\,u_k + K\,e_k,
\qquad
y_k = C\,x_k + D\,u_k + e_k.
\]

defined by interpreting $e_k = y_k-\hat y_{k|k-1}$ as innovation.

\pause
\medskip
For monic ARMAX, a convenient choice, consistent with the controllable
realization, is
\[
A := A_c,\quad B := B_c,\quad C := (C_u+C_e),\quad D := D_u,\quad K := B_c,\quad D_e := 1,
\]
so that
\vspace{-.5cm}
\begin{align*}    
    x_{k+1} &= A_c\,x_k + B_c\,u_k + B_c\,e_k,\\
    y_k &= (C_u+C_e)\,x_k + D_u\,u_k + e_k.
\end{align*}

% equivalent to the controllable form.

\pause
\medskip
\begin{alertblock}{}
\textbf{Innovation form}: the noise enters as a \emph{one-step-ahead prediction error}.
\end{alertblock}

\pause
\begin{alertblock}{}
Note that in the above SS model there is some abuse of notation since to
interpret $e_k$ as the innovation the model does not define the evolution of the
\tq{true} state $x_k$ but rather of the prediction $\hat{x}_k$.
\end{alertblock}

}%

%--------------------------
\myframe{Realization: Observer (Predictor) Form}{%

Assuming an innovation-form model
\[
x_{k+1} = A\,x_k + B\,u_k + K\,e_k,\qquad
y_k = C\,x_k + D\,u_k + e_k.
\]
Eliminate the unmeasured innovation $e_k$ using $e_k = y_k - (C\,x_k + D\,u_k)$.

\medskip
Substitute into the state update to get the \textbf{observer / predictor form}:
\[
x_{k+1} = (A-KC)\,x_k + (B-KD)\,u_k + K\,y_k,
\qquad
\hat y_{k|k-1} = C\,x_k + D\,u_k.
\]

\medskip
For consistency with the controllable ARMAX realization:
\[
A:=A_c,\quad B:=B_c,\quad C:=(C_u+C_e),\quad D:=D_u,\quad K:=B_c,
\]

\pause
\medskip
\begin{alertblock}{}
\textbf{Observer (predictor) form}: state propagated using the measured
output $y_k$.
\end{alertblock}

\medskip
\textbf{Note:} this form is not really generative. Rather $y_k$ are assumed to
be available.

}%


%--------------------------
\myframe{Realization: Continuous to Discrete (Exact) (part 1)}{%

Consider a general CT model
%
\begin{align*}
    \dot x(t) &= A\,x(t)+B\,u(t)+G\,w(t)\\
    y(t) &= C\,x(t)+D\,u(t)+v(t)
\end{align*}
%

\pause
Assume a sampling period $T>0$ and \textbf{zero-order hold (ZOH)} input, the
equations for $t\in[kT,(k+1)T]$ leads to
\[
x(t)=e^{A(t-kT)}x_k+\int_{kT}^{t} e^{A(t-\tau)}B\,u(\tau)\,d\tau
      +\int_{kT}^{t} e^{A(t-\tau)}G\,w(\tau)\,d\tau.
\]

\pause
and evaluate at $t=(k+1)T$ to
\[
x_{k+1}=e^{AT}x_k+\left(\int_{0}^{T} e^{A(T-\tau)}B\,d\tau\right)u_k
        +\int_{0}^{T} e^{A(T-\tau)}G\,w(kT+\tau)\,d\tau.
\]
}

\myframe{Realization: Continuous to Discrete (Exact) (part 2)}{

With the change of variable $\sigma=T-\tau$ and output sampling, we get:
\begin{align*}
    x_{k+1} &= A_d x_k + B_d u_k + w_k\\
    y_k &= C x_k + D u_k + v_k.
\end{align*}
%
\[
A_d:=e^{AT},\quad
B_d:=\int_{0}^{T} e^{A\sigma}B\,d\sigma,\quad
w_k:=\int_{0}^{T} e^{A\sigma}G\,w(kT+T-\sigma)\,d\sigma.
\]
%

\pause
If $w(t)$ is CT white noise with spectral density $Q_c$, $E[w(t)w(\tau)^\top]=Q_c\,\delta(t-\tau)$,
then $w_k$ is discrete-time white with
\[
Q_d := E[w_k w_k^\top] = \int_{0}^{T} e^{A\sigma}\,G\,Q_c\,G^\top\,e^{A^\top\sigma}\,d\sigma.
\]

\textbf{Note:} there are also other (simpler but inexact) discretization
techniques, e.g. Forward/Backward Euler etc.

}


%--------------------------
\myframe{Example of non-stationary process}{%

A process is \textbf{stationary} when its statistical properties (mean, variance,
autocorrelation) do not change over time. Typically arise from \textbf{stable systems}.

\begin{itemize}
    \item thermal noise in an electrical resistor
    \item vibrations in mechanical systems
    \item turbulent flow velocity fluctuations
    \item noise in sensors
\end{itemize}

\pause
\medskip
A process is \textbf{non-stationary} when its statistical properties change over time.
Often caused by: integration, accumulation, system drift.

\begin{itemize}
    \item particle under Brownian motion (diffusion process)
    \item position of a vehicle (tracking navigation)
    \item financial asset prices
\end{itemize}
}

%--------------------------
\myframe{State-Space Derivation of ARIMA$(p,d,q)$ (part 1)}{

Let $\Delta := 1-z^{-1}$ (difference operator). An ARIMA$(p,d,q)$ process satisfies
\[
A(z^{-1})\,\Delta^{d} y_k = C(z^{-1})\,e_k,\qquad e_k\sim WN(0,\sigma^2),
\]
with monic polynomials
\[
A(z^{-1}) = 1+a_1z^{-1}+\cdots+a_p z^{-p},\qquad
C(z^{-1}) = 1+c_1z^{-1}+\cdots+c_q z^{-q}.
\]

\pause
For an ARIMA model, define the stationary differenced series
\[
z_k := \Delta^{d} y_k.
\]
Then
\[
A(q^{-1})\,z_k = C(q^{-1})\,e_k
\]
is an ARMA$(p,q)$ model for $z_k$.
}

%--------------------------
\myframe{State-Space Derivation of ARIMA$(p,d,q)$ (part 2)}{
Find a $n$-dim innovation realization ($n=\max\{p,q+1\}$) for the ARMA for $z_k$
\[
x^z_{k+1} = A_z\,x^z_k + K_z\,e_k,\qquad
z_k = C_z\,x^z_k + e_k,
\]
where $(A_z,C_z)$ are chosen so that the transfer $e\mapsto z$ equals $\frac{C(q^{-1})}{A(q^{-1})}$.

\pause
\medskip
Add $d$ discrete integrators $s^{(1)}_k,\dots,s^{(d)}_k$ to recover $y_k=s^{(d)}_k$ from
$z_k=\Delta^d y_k$
\[
s^{(1)}_{k+1} = s^{(1)}_k + z_{k+1},\quad \dots,\quad
s^{(d)}_{k+1} = s^{(d)}_k + s^{(d-1)}_{k+1},
\]

Equivalently, with $s_k := \begin{bmatrix}s^{(1)}_k&\cdots&s^{(d)}_k\end{bmatrix}^\top$ and
$T_d$ the \tq{cumulative-sum} matrix:
%
\[
s_{k+1} = T_d\,s_k + b_d\,z_{k+1},\quad y_k = e_d^\top s_k,\quad b_d^\top=e_d^\top=\begin{bmatrix}0, \dots, 0, 1\end{bmatrix}
\]

\pause
\medskip
Let be $x_k := \begin{bmatrix}(x^z_k)^\top & s_k^\top\end{bmatrix}^\top$,
using the ARMA equations
%
\begin{align*}
    x_{k+1} &= 
        \begin{bmatrix}
            A_z & 0\\
            b_d C_z A_z & T_d
        \end{bmatrix}
        x_k
        +
        \begin{bmatrix}
            K_z\\
            b_d C_z K_z
        \end{bmatrix}
        e_k
        +
        \begin{bmatrix}
            0\\
            b_d
        \end{bmatrix}
        e_{k+1},\\
    y_k &= \begin{bmatrix}0 & e_d^\top\end{bmatrix} x_k.
\end{align*}

}

%--------------------------
\myframe{Example: tracking in 2D}{

Assume random acceleration driving the motion. The particle position motion is
\begin{align*}
\p_{k+1} &= \p_k + T\v_{k}, \qquad \p_k= \begin{bmatrix} x_k \\ y_k \end{bmatrix}\\
\v_{k+1} &= \v_{k} + T\a_{k},\qquad \a_{k}\sim WN(\mathbf{0}, \sigma^2_aI_2)
\end{align*}

\pause
Removing $\v_k$, the position satisfies the ARIMA$(0,2,0)$ model
\begin{align*}
    (1 - z^{-1})^2 \p_k &= \mathbf{p}_k - 2\mathbf{p}_{k-1} + \mathbf{p}_{k-2} = \e_k,
    \quad
    \mathbf{e}_k \sim WN(\mathbf{0},Q),\quad Q=T^2\sigma^2_aI_2\\
    \y_k &= \p_k + \w_k,\quad\w_k \sim WN(\mathbf{0},R)\quad \text{(measurements)}.
\end{align*}

\pause
\medskip
Following the above describe realization procedure, the state space model reads

\begin{align*}
    \x_{k+1} &=
    \begin{bmatrix}
        1 & 0 & T & 0 \\
        0 & 1 & 0 & T \\
        0 & 0 & 1 & 0 \\
        0 & 0 & 0 & 1
    \end{bmatrix}
    \x_k
    +
    \begin{bmatrix}
        0 & 0 \\
        0 & 0 \\
        1 & 0 \\
        0 & 1
    \end{bmatrix}
    \a_k,
    \quad
    \a_k \sim WN(\mathbf{0},\sigma^2_aI_2).\\
    %
    \y_k &=
        \begin{bmatrix}
            1 & 0 & 0 & 0 \\
            0 & 1 & 0 & 0
        \end{bmatrix}
    \x_k + \w_k\qquad\qquad\qquad \w_k \sim WN(\mathbf{0},R).
\end{align*}

}

%-------------------------------------------------------------------------------
\subsection{The Kalman Filter(s) and Predictor}

%--------------------------
\myframe{The Kalman Filter: reference state space model}{

In the following we consider the general state space model of the process
$\{\y(t)\}$

\begin{align*}
    \x(t+1) &= A\x(t) + \v(t)\\
    \y(t) &= C\x(t) + \w(t)\quad
    \E\begin{bmatrix}\v(t)\\\w(t)\end{bmatrix} = 0,\
    \Var\begin{bmatrix}\v(t)\\\w(t)\end{bmatrix} = \begin{bmatrix}Q & S \\ S^T & R \end{bmatrix}, R\succ0 \\
    \x(t_0) &= \xzero,\quad \E[\xzero] = \muzero,\quad \Var[\x_0]=\Pzero\succcurlyeq 0,\quad \xzero\ \bot\ \v(t),\w(t)\ \forall t\geq t_0
\end{align*}

\pause
\vspace{.5cm}
\begin{tabular}{p{7cm}cc}
As \textit{widely} known, the Kalman filter deals with
estimation of the (unknown and usually not observable) state process $\{\x(t)\}$ &
\pause
$\Longrightarrow$ &
\textbf{But why?}
\end{tabular}

\pause
\vspace{.5cm}
\textbf{Note}: $A,C,Q,S,R$ could be time varying. $R\succ 0$ for simplicity for
now because it helps guaranteeing invertibility of the variance of innovation
process $\e(t)=\y(t)-\hat{\y}(t|t-1)$.

}

%--------------------------
\myframe{The Kalman Filter: why the state $\x(t)$? (part 1)}{

Two prototypical problems:

\pause
\vspace{.3cm}
\textbf{1. Prediction} of $\{\y(t)\}$ from past measurements.

\pause
\medskip
By observing that $\w(t+1)\ \bot\ \{\v(s),\w(s),\xzero,s\leq t\}\supset\Hpast{t}(\y)$
%
\pause
$$
\Longrightarrow \yhat(t+1|t) = \Elin[\y(t+1)|\Hpast{t}(\y)] = C\Elin[\x(t+1)|\Hpast{t}(\y)]=C\xhat(t+1|t)
$$

\pause
\textbf{The prediction can be directly built from the predicted state!}

\pause
\vspace{.3cm}
\textbf{2. Filtering} of $\{\z(t)\}$ from $\{\y(t)\}$

\pause
\medskip
Say we want to filter $\z(t)$ from $\y(t)=\z(t)+\xiv(t)$ ($\xiv$ \tq{colored noise}).
\pause
Assume the following SS representation for the processes:
%
\pause
\begin{align*}
    \x(t+1) &=
        \begin{bmatrix}
            \x_1(t+1) \\ \x_2(t+1)
        \end{bmatrix}
        =
        \begin{bmatrix}
            A_1 & \\ & A_2
        \end{bmatrix}
        \begin{bmatrix}
            \x_1(t) \\ \x_2(t)
        \end{bmatrix}
        +
        \begin{bmatrix}
            \v_1(t) \\ \v_2(t)
        \end{bmatrix}
        = A\x(t) + \v(t)\\
    \y(t) &=
        \z(t) + \xiv(t) =
        \begin{bmatrix} C_1 & C_2 \end{bmatrix}
        \begin{bmatrix} \x_1(t) \\ \x_2(t) \end{bmatrix}
        +
        \begin{bmatrix} I & I\end{bmatrix}
        \begin{bmatrix} \w_1(t) \\ \w_2(t)\end{bmatrix}
        =
        C\x(t) + \w(t)\\
    \z(t) &=
        \begin{bmatrix}C_1 & 0\end{bmatrix}
        \begin{bmatrix} \x_1(t) \\ \x_2(t) \end{bmatrix}
        +
        \w_1(t)
        = \overline{C}_1\x(t) + \w_1(t)
\end{align*}

}

%--------------------------
\myframe{The Kalman Filter: why the state $\x(t)$? (part 2)}{
$$
\Longrightarrow \zhat(t|t)=\Elin[\z(t)|\Hpast{t}(\y)] = \overline{C}_1\xhat(t|t) + \what_1(t|t)
$$

\pause
\vspace{.3cm}
But what about $\what_1(t|t)$?
%
$$
\w(t)\ \not\perp\ \w_1(t)\ \textnormal{with}\ \E[\w_1(t)\w(t)^\top]=R_1 \Longrightarrow \y(t)\ \not\perp\ \w_1(t)
$$
%

\pause
Nonetheless, via double projection (on subspaces)
%
\pause
\begin{align*}
    \what_1(t|t) &= \Elin\left[\Elin\left[\w_1(t)| \{\xzero, \v(s-1), \w(s); s\leq t\}\right] | \Hpast{t}(\y)\right]\\
    &= \Elin\left[R_1R^{-1}\w(t) | \Hpast{t}(\y)\right]\\
    &= R_1R^{-1}\what(t|t)\\ 
    &= R_1R^{-1}\left(\y(t) - C\xhat(t|t)\right)
\end{align*}

\vspace{-.5cm}
\pause
\[
    \Longrightarrow\quad  \zhat(t|t) = \overline{C}_1\xhat(t|t) + R_1R^{-1}\left(\y(t) - C\xhat(t|t)\right)
\]

\pause
\begin{alertblock}{}
    \textbf{Again, the filtering solution can be directly built for the filtered
    state!}
\end{alertblock}
}


%--------------------------
\myframe{The Kalman Filter: reference state space model (modified)}{

\begin{align*}
    \x(t+1) &= A\x(t) + \v(t)\\
    \y(t) &= C\x(t) + \w(t)\quad
    \E\begin{bmatrix}\v(t)\\\w(t)\end{bmatrix} = 0,\
    \Var\begin{bmatrix}\v(t)\\\w(t)\end{bmatrix} = \begin{bmatrix}Q & S \\ S^T & R \end{bmatrix}, R\succ0 \\
    \x(t_0) &= \xzero,\quad \E[\xzero] = \muzero,\quad \Var[\x_0]=\Pzero\succcurlyeq 0,\quad \xzero\ \bot\ \v(t),\w(t)\ \forall t\geq t_0
\end{align*}

\pause
\medskip
Note that $\E[\v(t),\w(t)^T]=S$ but $\v(s)\perp\w(t),\ s\neq t$.

\pause
\begin{align*}
    \Longrightarrow \verr(t) = \v(t) - \Elin[\v(t)|\H(\w)] &= \v(t) - \Elin[\v(t)|\w(t)]\qquad \perp \w(t)\\
    &= \v(t) - SR^{-1}\w(t)\\
    \Longrightarrow \Var\begin{bmatrix} \verr(t)\\ \w(t)\end{bmatrix} = \begin{bmatrix}\tilde{Q} &  \\  & R \end{bmatrix},\quad\,\,\, & \tilde{Q}=Q-SR^{-1}S^T
\end{align*}

\pause
$$
    \Longrightarrow \x(t+1) = A\x(t) + \v(t) = F\x(t) + SR^{-1}\y(t) + \verr(t),\quad F=A-SR^{-1}C
$$

with uncorrelated noise $\verr(t)\perp\w(t)$.
}


%--------------------------
\myframe{The Kalman Filter: equations}{%
    \begin{enumerate}
        \pause
        \item A-priori estimates/predictions (temporal update) 
        \begin{itemize}
            \item State: $\xhat(t+1|t) = F \xhat(t|t) + SR^{-1}\y(t)$
            \item Covariance: $P(t+1|t) = F P(t|t) F^T + \tilde{Q}$, $\quad t\geq t_0$
        \end{itemize}

        \pause
        \medskip
        \item A-posteriori estimates/filters (measurement update)
        \begin{itemize}
            \item State: $\xhat(t+1|t+1) = \xhat(t+1|t) + L(t+1) (\y(t+1) - C \xhat(t+1|t))$
            \item Covariance: $P(t+1|t+1) = P(t+1|t) - P(t+1|t)C^T\Lambda(t+1)^{-1}CP(t+1|t)$
        \end{itemize}

        \pause
        \medskip
        \item Initial conditions
        \begin{itemize}
            \item $\xhat(t_0|t_0-1)=\muzero,\qquad P(t_0|t_0-1)=\Pzero$
        \end{itemize}
    \end{enumerate}

    \pause
    where
    
    \begin{itemize}
        \item $P(t+1|t)$, $P(t|t)$: the prediction and filtering \underline{\textbf{errors}} variances
        \item $\Lambda(t)=CP(t|t-1)C^T + R$: innovation $(\e(t)=\y(t)-C\xhat(t|t-1))$ variance
        \item $L(t) = P(t|t-1)C^T\Lambda(t)^{-1}$: Kalman gain
    \end{itemize}
}


%--------------------------
\myframe{The Kalman Filter: equations - proof (part 1)}{%
% To obtain the recursions given in the previous slide we exploit the
% conditional-orthogonality properties of the state-space model together with the
% representation

% $$
%     \x(t+1) = F\x(t) + SR^{-1}\y(t) + \verr(t),
%     \qquad \verr(t)\perp \w(s) \ \forall s \le t
% $$

1. \textbf{A-priori (prediction) update.}  The one-step ahead estimate is
obtained by taking the linear conditional expectation given past measurements
$\Hpast{t}(\y)$, recalling that $\verr(t)\perp\w(t)\Longrightarrow\verr(t)\perp\Hpast{t}(\y)$:
%
$$
    \xhat(t+1|t) := \Elin[\x(t+1)|\Hpast{t}(\y)]
            = F\Elin[\x(t)|\Hpast{t}(\y)] + SR^{-1}\y(t)
            = F\xhat(t|t) + SR^{-1}\y(t)
$$

The corresponding prediction error
$\xerr(t+1|t)=\x(t+1)-\xhat(t+1|t)$ satisfies

$$
    \xerr(t+1|t)=F\xerr(t|t)+\verr(t)
$$

and, since $\verr(t)\perp\xerr(t|t)$ because
%
$$
\left.
\begin{aligned}
    &\x(t)\in\spn{\x(t_0), \verr(s), \w(s);\ s<t}\\
    &\Hpast{t}(\y)\subset \spn{\x(t_0), \verr(s-1), \w(s); s\leq t}
\end{aligned}
\right\}\perp\verr(t)
$$

its covariance evolves as
%
$$
    P(t+1|t) = \Var[\xerr(t+1|t)]
             = F P(t|t) F^{T} + \tilde Q
$$
}

\myframe{The Kalman Filter: equations - proof (part 2)}{

2. \textbf{A-posteriori (measurement) update.}  At time $t+1$ we get
$\y(t+1)=C\x(t+1)+\w(t+1)$. Note that
%
$$
\Hpast{t+1}(\y) = \Hpast{t}(\y) \oplus \H(\e(t+1))\quad \textnormal{(orthogonal sum)}
$$
%
being
$$
\e(t+1) = \y(t+1) - C\xhat(t+1|t) = C\xerr(t+1|t) + \w(t+1)
$$
\vspace{-.2cm}
\begin{align*}
    \Longrightarrow \xhat(t+1|t+1) &= \E[x(t+1)|\Hpast{t+1}(\y)] = \xhat(t+1|t) + L(t+1)\e(t+1)\\
    &\xhat(t+1|t) + L(t+1)\big(\y(t+1)-C\xhat(t+1|t)\big)
\end{align*}

\medskip
where $L(t+1)=\Cov[\x(t+1),\e(t+1)]\Var[\e(t+1)]^{-1}$ with
\vspace{.1cm}
\begin{align*}
    &\Cov[\x(t+1),\e(t+1)] = \Cov[\x(t+1),\xerr(t+1|t)]C^T + \Cov[\x(t+1),\w(t+1)]\\
    &\qquad\qquad\qquad\qquad\quad\ \ = \Cov[\xerr(t+1|t),\xerr(t+1|t)] = P(t+1|t)\\ 
    &\Var[\e(t+1)] = \Lambda(t+1) = CP(t+1|t)C^T + R,\quad \textnormal{since}\ \xerr(t+1|t)\perp\w(t+1)  
\end{align*}

}

\myframe{The Kalman Filter: equations - proof (part 3)}{

The covariance update follows by noting that

\begin{align*}
    &\x(t+1) - \xhat(t+1|t+1) = \x(t+1) - \xhat(t+1|t) - L(t+1)\e(t+1)\\
    &\Longrightarrow \xerr(t+1|t+1) = \xerr(t+1|t) - L(t+1)\e(t+1)\\
    &\Longrightarrow \xerr(t+1|t+1) + L(t+1)\e(t+1) = \xerr(t+1|t) 
\end{align*}

and being $\xerr(t+1|t+1) \perp \e(t+1)$

$$
    P(t+1|t+1) + L(t+1)\Lambda(t+1)L(t+1)^T = P(t+1|t)
$$
%
\begin{align*}
    \Longrightarrow P(t+1|t+1) &=  P(t+1|t) - L(t+1)\Lambda(t+1)L(t+1)^T\\
                               &= P(t+1|t) - P(t+1|t)C^T\Lambda(t+1)^{-1}CP(t+1|t)
\end{align*}

}

%--------------------------
\myframe{Considerations on error covariance update}{%

\begin{enumerate}
    \item To avoid/reduce numerical ill-conditioning it is sometimes preferred
    to express the covariance update in different forms
    \begin{align*}
        P(t+1|t+1) &=  P(t+1|t) - P(t+1|t)C^T\Lambda(t+1)^{-1}CP(t+1|t)\\
                   &= \left( I - L(t+1)C\right) P(t+1|t) \\
       \textnormal{symmetrized form} &= \left[I-L(t+1)C\right] P(t+1|t) \left[I-L(t+1)C\right]^T +\\
                   &\quad L(t+1)RL(t+1)^T
    \end{align*}
    
    \pause
    \item Note that \textbf{ALL} the \textbf{covariance computations} (error and
    innovation) as well as the \textbf{Kalman Gain} do not explicitly depend on the
    measurements and can be pre-computed \textbf{offline}. This is the major computational burden!
    
    \pause
    \medskip
    \item Conversely, the \textbf{state update and prediction} depend on the measurements
    and must be computed \textbf{online}.
\end{enumerate}


}


%--------------------------
\myframe{The Kalman Predictor}{%

It is often useful to express the dynamics of the predictor which allows to
express the next state prediction as function of the previous without computing
the filtered value. That is (after some manipulation)
%
\pause
\begin{align*}
    \xhat(t+1|t) &= F(I - L(t)C)\xhat(t|t-1) + (FL(t) + SR^{-1})\y(t)\\
                 &= (F - K(t)C)\xhat(t|t-1) + (K(t) + SR^{-1})\y(t)\\
                 &= \Gamma(t)\xhat(t|t-1) + G(t)\y(t)\\
                 &= A\xhat(t|t-1) + G(t)\left[\y(t)-C\xhat(t|t-1)\right]\\
                 &= A\xhat(t|t-1) + G(t)\e(t)\\
    P(t+1|t) &= F \left[P(t|t-1) - P(t|t-1)C^T\Lambda(t)^{-1}CP(t|t-1)\right] F^T + \tilde{Q}\\
             &= F\left[I-L(t)C\right]P(t|t-1)\left[I-L(t)C\right]^TF^T + FL(t)RL^T(t)F^T + \tilde{Q}
\end{align*}
%
where $K(t):=FL(t)$, $\Gamma(t):= F-K(t)C$, $G(t):=K(t)+SR^{-1}$.
}

%--------------------------
\myframe{Equivalent closed loop interpretation}{
% \centering
\begin{tikzpicture}[
    auto,
    >=Latex,
    node distance=1.8cm and 2.2cm,
    block/.style={draw, rectangle, rounded corners, minimum width=1cm, minimum height=1.1cm, align=center},
    sum/.style={draw, circle, inner sep=1pt, minimum size=6mm},
    input/.style={align=center},
    every path/.style={draw, -Latex}
]

% Nodes
\node[input] (y) {$\y(t)$};
\node[sum, right=.8cm of y] (innovsum) {$-$};
\node[block, right=1.0cm of innovsum] (gain) {$G(t)$};
\node[sum, right=.8cm of gain] (statesum) {$+$};
\node[block, right=.8cm of statesum] (z) {$z^{-1}$};
\node[input, right=1.5cm of z] (x) {$\xhat(t|t-1)$};

\node[block, below=1.0cm of z] (A) {$A$};
\node[block, below=2.0cm of gain] (C) {$C$};


% Connections
\path (y) -- (innovsum);
\path (innovsum) -- node {$\e(t)$} (gain);
\path (gain) -- (statesum);
\path (statesum) -- (z);
\path (C) -| node[pos=0.25, right] {$-$} (innovsum);
\path (A) -| node[pos=0.25, right] {} (statesum);
\path (x) |- node[pos=0.25, right] {} (A);
\path (x) |- node[pos=0.25, right] {} (C);
\path (z) -- (x);

\end{tikzpicture}

where the closed-loop dynamics is described by
$$
\Gamma(t) = A-G(t)C = F - K(t)C
$$
Hence its properties characterize the filter/predictor dynamics.
}

%--------------------------
\myframe{Filtered Estimates}{%

Similarly (for $S=0$ for simplicity)

\begin{align*}
    \xhat(t+1|t+1) &= A\xhat(t|t) + L(t+1)\left[\y(t+1)-CA\xhat(t|t)\right]\\
     P(t+1|t+1) &= \left[I-L(t+1)C\right]P(t+1|t)
\end{align*}

In this case the closed-loop dynamical behavior is described by

$$
\Gamma_{\rm filt}(t) = \left[I - L(t+1)C\right]A
$$

where before, with $S=0$, was $\Gamma_{\rm pred}(t)=A\left[I - L(t)C\right]$ which suggest
an almost identical asymptotic dynamic behavior.

\vspace{1cm}

\pause
\textit{Note that prediction and filtered estimates can directly be computed
from the original equations. It is sometimes convenient to get the explicit
expressions in order to better understand the role of the measurements in the
update.}

}

%--------------------------
\myframe{Kalman Filter with Input (part 1)}{
Reference model (with usual assumptions on noise and initial state):
%
\begin{align*}
    \x(t+1) &= A\x(t) + B\u(t) + \v(t) = F\x(t) + SR^{-1}\y(t) + B\u(t) + \verr(t)\\
    \y(t) &= C\x(t) + \w(t) + D\u(t)
\end{align*}


\pause
\textbf{Easy}: if $\u(t)$ \textbf{known and deterministic}
%
\begin{align*}
    \xhat(t+1|t) &= F\xhat(t|t) + SR^{-1}\y(t) + B\u(t)\\
    P(t+1|t) &= ...\ unchanged\\
    \e(t) &= y(t) - C\xhat(t|t-1) - D\u(t)
\end{align*}
%
$\Longrightarrow$ the deterministic input does not affect the uncertainty, only
the trajectory.

\pause
\medskip
\textbf{Harder}: \textit{but what if $\u(t)$ is a process?} E.g. a (linear) causal function ($H$) of
$\Hpast{t}(\y)$ plus a process uncorrelated from the entire noise ($\v, \w$)
history, say

$$
\u(t) = H(\y^t) + \r(t)
$$

\pause
In this case, even if $H\equiv 0$, it is in general not possible to assume
$\u(t)\perp\xzero$. Even if $\u(t)$ is observable, $\xhat$ based
on joint observations $(\y, \u)$ must be obtained by considering the joint state
model thus deriving augmented KF equations.

}

%--------------------------
\myframe{Kalman Filter with Input (part 2)}{

Assume the plant and input are described by

\begin{align*}
    &\begin{cases}
        \x_1(t+1) &= A\x_1(t) + B\u(t) + \v_1(t)\\
        \y(t) &= C \x(t) + \w_1(t)
    \end{cases}\\
    %
    &\begin{cases}
        \x_2(t+1) &= F\x_2(t) + G\y(t) + \v_2(t)\\
        \u(t) &= H\x_2(t) + J\y(t) + \w_2(t)
    \end{cases}
\end{align*}

\pause
This gives rise to the augmented state model to which standard KF can be applied

\begin{align*}
    \x(t+1) &=
    \underbrace{
        \begin{bmatrix}
            A+BJC & BH\\
            GC & F
        \end{bmatrix}
    }_{\mathcal A}
    \x(t)
    +
    \underbrace{
        \begin{bmatrix}
            \v_1(t) \\ \v_2(t)
        \end{bmatrix} +
        \begin{bmatrix}
            BJ & B\\
            G & 
        \end{bmatrix}
        \begin{bmatrix}
            \w_1(t) \\ \w_2(t)
        \end{bmatrix}
    }_{\overline{\v}(t)}\\
%
    \begin{bmatrix}
        \y(t)\\ \u(t)
    \end{bmatrix}
    &=
    \underbrace{
        \begin{bmatrix}
            C &  \\
            JC & H
        \end{bmatrix}
    }_{\mathcal C}
    \x(t) + 
    \underbrace{
        \begin{bmatrix}
            I & \\ J & I
        \end{bmatrix}
    }_{\mathcal{D}}
    \underbrace{
        \begin{bmatrix}
            \w_1(t) \\ \w_2(t)
        \end{bmatrix}
    }_{\w(t)}
\end{align*}

}

%--------------------------
\myframe{The Kalman Filter: information form}{%
    \begin{enumerate}
        \item Information variables
        \begin{itemize}
            \item Information matrix: $Y(t|t) \triangleq P(t|t)^{-1}$
            \item Information state: $\xi(t|t) \triangleq Y(t|t)\xhat(t|t)$
        \end{itemize}

        \pause
        \medskip
        \item A-priori estimates/predictions (temporal update)
        \begin{itemize}
            \item State:
                $
                \xi(t+1|t)
                =
                Y(t+1|t)
                \left(
                    F\,Y(t|t)^{-1}\xi(t|t)
                    +
                    SR^{-1}\y(t)
                \right)
                $
            \item Information matrix:
                $
                Y(t+1|t)
                =
                \left(
                    F\,Y(t|t)^{-1}F^T + \tilde{Q}
                \right)^{-1},
                \quad t\geq t_0
                $
        \end{itemize}

        \pause
        \medskip
        \item A-posteriori estimates/filters (measurement update)
        \begin{itemize}
            \item Information state:
                $
                \xi(t+1|t+1)
                =
                \xi(t+1|t) + C^T R^{-1}\y(t+1)
                $
            \item Information matrix:
                $
                Y(t+1|t+1)
                =
                Y(t+1|t) + C^T R^{-1} C
                $
        \end{itemize}

        \pause
        \medskip
        \item Initial conditions
        \begin{itemize}
            \item $Y(t_0|t_0-1)=\Pzero^{-1},\qquad \xi(t_0|t_0-1)=\Pzero^{-1}\muzero$
        \end{itemize}
    \end{enumerate}

    \pause
    where
    
    \begin{itemize}
        \item $Y(t|t-1)$, $Y(t|t)$: the prediction and filtering \underline{\textbf{information}} matrices
        \item $\xi(t|t-1)$, $\xi(t|t)$: the corresponding information states
        \item $\xhat(t|t)=Y(t|t)^{-1}\xi(t|t)$
    \end{itemize}
}

%--------------------------
\myframe{Kalman filter vs. Information filter}{%
    \begin{columns}[t]
        \column{0.49\textwidth}
        \textbf{Standard Kalman filter}
        \begin{itemize}
            \item Propagates
            \[
                \xhat(t|t),\quad P(t|t)
            \]
            \item Measurement update through
            \[
                \Lambda(t)=CP(t|t-1)C^T+R
            \]
            \item Gain computation required:
            \[
                L(t)=P(t|t-1)C^T\Lambda(t)^{-1}
            \]
            \item Natural for moderate state dimension and direct state reconstruction at each $t$
            \item Direct Covariance update
        \end{itemize}

        \column{0.49\textwidth}
        \textbf{Information filter}
        \begin{itemize}
            \item Propagates
            \[
                Y(t|t),\quad \xi(t|t)
            \]
            \item Measurement update is additive:
            \[
                Y^+ = Y^- + C^TR^{-1}C
            \]
            \[
                \xi^+ = \xi^- + C^TR^{-1}y
            \]
            \item No explicit Kalman gain
            \item Natural for multi-sensor fusion and sparse large-scale problems
            \item Particularly convenient when measurements are many and independent
        \end{itemize}
    \end{columns}

    \medskip

    \textbf{Key relation}
    \[
        Y(t|t)=P(t|t)^{-1},
        \qquad
        \xi(t|t)=Y(t|t)\xhat(t|t)
    \]
}

%--------------------------
\myframe{Why information form can be advantageous}{%
    \textbf{Example: fusion of many independent sensors}

    Assume $N$ sensors observe the same state:
    \[
        y_i(t)=C_i x(t) + v_i(t),
        \qquad
        v_i(t)\sim \mathcal N(0,R_i),
        \qquad i=1,\dots,N
    \]

    \pause
    \medskip
    \textbf{Standard Kalman filter}
    \begin{itemize}
        \item Meas. typically processed one by one, or stacked into one large vector
        \item This leads to repeated gain computations and innovation covariances
    \end{itemize}

    \pause
    \medskip
    \textbf{Information filter}
    \begin{itemize}
        \item Sensor contributes an additive information term (computed independently):
        \[
            Y(t|t)
            =
            Y(t|t-1)
            +
            \sum_{i=1}^{N} C_i^T R_i^{-1} C_i;
        \quad
            \xi(t|t)
            =
            \xi(t|t-1)
            +
            \sum_{i=1}^{N} C_i^T R_i^{-1} y_i(t)
        \]
        \item Very convenient for
        \begin{itemize}
            \item distributed sensor networks
            \item sparse estimation problems such as SLAM
        \end{itemize}
    \end{itemize}

    \pause
    \medskip
    \textbf{Main advantage}
    
    Updates are \underline{\textbf{simple accumulation of
information}}: more scalable and simpler.}



%--------------------------------
\subsection{Optimality and Asymptotic Theory}

%--------------------------------
\myframe{Cramer-Rao Lower Bound (CRLB) for KF}{
    
    \pause
    The CRLB defines the minimum theoretical variance of \underline{any
    unbiased estimator} of a \underline{deterministic parameter}, setting a
    fundamental limit on estimation precision. It is defined as the inverse of
    the Fisher information.

    \pause
    \medskip
    In the KF context, we deal with estimation of a \underline{time-varying
    random vector based on data and prior information}. Hence the relevant bound
    is what is usually referred to as the \textbf{posterior (Bayesian) CRLB}
    %
    $$
        \mathcal{I}(t|t) = -\E[\nabla_\x^2 \log p(\x(t)|\Hpast{t}(\y))\ |\ \Hpast{t}(\y)]\qquad \textnormal{Fisher Information Matrix}
    $$
    
    \medskip
    \begin{itemize}
        \pause
        \item For any unbiased estimator $\xhat_u(t|t)$ of $\x(t)$ from $\Hpast{t}(\y)$:
        \[
            \E[\xhat_u(t|t)-\x(t)] = 0,
            \qquad
            \Cov[\xhat_u(t|t)-\x(t)] \succcurlyeq \mathcal{I}(t|t)^{-1}
        \]
        %where $\mathcal{I}(t|t)$ is the Fisher information matrix.

        \pause
        \medskip
        \item Under the linear-Gaussian assumptions used for KF:
        \[
            \mathcal{I}(t|t)=Y(t|t)=P(t|t)^{-1} \quad \Longrightarrow \quad \Cov[\xhat_u(t|t)-\x(t)] \succcurlyeq P(t|t)
        \]
        and the Kalman filter attains the bound:
        \[
            \Cov[\xerr(t|t)]
            =
            \Cov[\x(t)-\xhat(t|t)]
            =
            P(t|t).
        \]
    \end{itemize}
}

%--------------------------
\myframe{Cramer-Rao limit (KF): proof of bound attainability}{
    Under the linear-Gaussian assumptions
    $\x(t)\,|\,\Hpast{t}(\y)\sim\mathcal N\!\left(\xhat(t|t),\,P(t|t)\right)$, where

    \[
        \xhat(t|t)=\Elin[\x(t)|\Hpast{t}(\y)]\equiv \E[\x(t)|\Hpast{t}(\y)]\quad (\textnormal{the KF})
    \]

    \medskip
    For any estimator $\xhat_u(t|t)$ based on $\Hpast{t}(\y)$, define
    \[
        d(t):=\xhat_u(t|t)-\xhat(t|t),
        \qquad
        \xerr(t|t):=\x(t)-\xhat(t|t).
    \]
    
    \medskip
    By the orthogonality principle of conditional expectation,
    % \[
    $\E\!\left[\xerr(t|t)\,d(t)^T\right]=0$.
    % \]

    \medskip
    Hence
    \begin{align*}
        \Cov[\x(t)-\xhat_u(t|t)]
        &= \Cov[\xerr(t|t)-d(t)]\\
        &= \Cov[\xerr(t|t)] + \Cov[d(t)]\\
        &= P(t|t) + \Cov[d(t)]
        \succcurlyeq P(t|t).
    \end{align*}

    Equality holds iff $\Cov[d(t)]=0$ (i.e., $d(t)=0$ a.s.), namely
    $\xhat_u(t|t)=\xhat(t|t)$
}

% %--------------------------
% \myframe{Asymptotic Behavior}{%

% \todo{add asymptotic behavior}

% }